\documentclass[12pt]{article}
\usepackage[a4paper, margin=2cm]{geometry}

\usepackage{fontspec}
\usepackage{polyglossia}
\setdefaultlanguage{russian}
\setmainlanguage{russian}
\setmainfont{Times New Roman}

\begin{document}

\begin{center}
\textbf{\LARGE ЗАЯВЛЕНИЕ}
\end{center}

\vspace{1cm}



% === ТЕКСТ ДОКУМЕНТА ВСТАВЛЯТЬ СЮДА ===
% Здесь должно быть изложено обращение, просьба или заявление
% в свободной форме, но в официально-деловом стиле.

Настоящим уведомляем, что по состоянию на дату \textbf{26 марта 2026 года} в предоставленном информационном контексте отсутствуют верифицируемые сведения о \textbf{текущих котировках акций} основных нефтеперерабатывающих компаний (далее — НПЗ-компании), включая их рыночные тикеры и соответствующие значения котировок на указанную дату.

По результатам проверки доступного контекста установлено, что присутствуют фрагментарные данные, относящиеся преимущественно к рынку сырой нефти (например, котировки/уровни Brent и WTI в отдельные моменты времени), однако \textbf{не содержится} перечня нефтеперерабатывающих компаний и их акций, по которым можно составить сравнительный отчет котировок именно на \textbf{26 марта 2026 года}.

Учитывая изложенное, формирование запрошенного отчета в части \textbf{“актуальных котировок акций основных нефтеперерабатывающих компаний на 26.03.2026”} в рамках текущего контекста не представляется возможным. Предлагается рассмотреть один из следующих вариантов:

\begin{enumerate}
\item предоставить перечень компаний (или тикеров) и/или исходные источники котировок (с указанием дата-временной отметки и биржи/площадки);
\item разрешить использование альтернативного набора источников (в том числе с финансовых терминалов/провайдеров данных), обеспечивающих подтверждаемость котировок на дату 26 марта 2026 года.
\end{enumerate}

Дополнительно сообщаем, что в предоставленном контексте встречались сведения, относящиеся к одной российской компании нефтепереработки \textbf{(KRKNP)}, однако указанные значения относились к \textbf{другой дате (27.03.2026)} и потому не могут быть использованы для расчета/сопоставления по состоянию на \textbf{26.03.2026}.

В случае предоставления перечня НПЗ-компаний и источников котировок обязуемся подготовить официальный отчет, включающий:
(а) список компаний и идентификаторы (тикеры);
(б) котировки акций на дату \textbf{26 марта 2026 года};
(в) при необходимости — пояснение методики (биржа/площадка, валюта, время фиксации, тип цены: last/close/settlement).

\vfill

\noindent
Дата: \rule{5cm}{0.4pt} \hfill Подпись: \rule{5cm}{0.4pt}

\end{document}